\subsection*{Auxiliary calculation: a nonsplit cross-factor refinement}

This calculation records a finite refinement of D.38 with a noncentral
kernel.  It gives no additional exclusion among the tested base tuples.
The complete systems, the specified cocycle, and their executed outputs
are retained so that this outcome can be checked independently.
No claim number is assigned to this auxiliary calculation.

\paragraph{The explicit nonsplit extension.}
Use the permutation numbering of D.38 and its checked identification
$A=A_5=\operatorname{PSL}_2(\F_5)$ from D.45.  Put
$V=\mathfrak{sl}_2(\F_5)$, with basis
\[
 \begin{pmatrix}1&0\\0&-1\end{pmatrix},\qquad
 \begin{pmatrix}0&1\\0&0\end{pmatrix},\qquad
 \begin{pmatrix}0&0\\1&0\end{pmatrix}.
\]
Write $\rho$ for conjugation on these matrices.  The program
\path{code/part_a/a5_cross_module.py} verifies the representation law on
all $3600$ pairs, and row reduction shows that the span of the sixty
matrices $\rho(g)$ is the full $9$-dimensional matrix algebra.
In particular $V$ is absolutely simple, $V^A=0$, and
$\operatorname{End}_A(V)=\F_5$.
The same program enumerates the crossed homomorphisms using the two
triangle generators with indices $3,19$ and obtains
\[
 \dim Z^1(A,V)=4,\qquad \dim B^1(A,V)=3,
                 \qquad \dim H^1(A,V)=1.
\]
It retains all sixty values of a representative $c$ of the nonzero
class, with $c(3)=(0,0,0)$ and $c(19)=(0,1,1)$, and verifies
$c(gh)=c(g)+\rho(g)c(h)$ for every pair.

For $A\times A$ acting on $V\otimes V$, define
\[
 b\bigl((g,h),(u,v)\bigr)=c(g)\otimes\rho(h)c(v).       \tag{Aux.1}
\]
The crossed-homomorphism identity gives the two-cocycle identity
directly.  For example, the two middle terms on either side agree
after expanding $c(gu)$ and $c(vw)$.  The program additionally checks
$57\,600$ generating-grid instances.  This cocycle is not a
coboundary.  Indeed, a cochain with coboundary $b$ would restrict to
crossed homomorphisms $a(g)$ and $d(h)$ on the two factors, since
$b$ vanishes on either individual factor.  Comparing the orders of
the two commuting factors gives
\[
 c(g)\otimes c(h)=(\rho(g)-1)d(h)-(\rho(h)-1)a(g).
\]
Taking first-factor cohomology yields
$[c]\otimes c(h)=-(\rho(h)-1)[a]$ in
$H^1(A,V)\otimes V$.  Writing $[a]=[c]\otimes w$ would make
$c(h)=-(\rho(h)-1)w$, contrary to the enumerated nonzero class.

Also $H^1(A\times A,V\otimes V)=0$.  The same commuting-factor
comparison for a one-cocycle shows that its restriction class to
the first factor lies in
$H^1(A,V)\otimes V^A=0$.  Subtract a coboundary to make that
restriction zero.  Its second-factor values are then fixed by the
first factor and are zero because $(V\otimes V)^{A\times1}=0$.
Tensoring the matrix spans proves that $V\otimes V$ is simple and
has scalar endomorphisms.

Take the following orbit of ordered factor pairs under the actual
$\theta,\omega$ component maps of D.38:
\[
 \mathcal O=\{(3,13),(5,15),(6,8),(7,8),(11,13),(14,15)\}.
\]
Its support is
$I=\{3,5,6,7,8,11,13,14,15\}$.  The projection of the specified
quotient $A_5^{19}$ to these factors is the specified epimorphism
$F\twoheadrightarrow Q=A^I$.  Put
\[
 M=\bigoplus_{(i,j)\in\mathcal O} V_i\otimes V_j,
 \qquad \dim_{\F_5}M=54,
 \qquad
 B(q,r)_{i,j}=c(q_i)\otimes\rho(q_j)c(r_j).
\]
The finite group used in the test is the set $M\times Q$ with
the exact multiplication
\[
 (m,q)(n,r)=(m+q n+B(q,r),qr).                            \tag{Aux.2}
\]
Denote it by $L$.  It has order $5^{54}60^9$.
The six summands of $M$ are pairwise nonisomorphic simple
$Q$-modules, as their two-factor supports are distinct.  Each
component of the extension class is nonzero by (Aux.1).
Consequently every subgroup of $L$ surjecting to $Q$ equals $L$:
its intersection with $M$ is a $Q$-submodule, and an omitted
simple summand would give a splitting of that component after
quotienting by the intersection.  Thus $M\subset\Phi(L)$.

The group $L$ is center-free: a central element projects to
$Z(Q)=1$ and then lies in $M^Q=0$.  It is also perfect, since
$Q$ is perfect and $[Q,M]=M$.  Finally every automorphism of $L$
over the identity of $Q$ is conjugation by an element of $M$.
Its module map is scalar on each of the six summands.  Preservation
of the nonzero extension component forces every scalar to be one.
The remaining difference is a one-cocycle, and
$H^1(Q,M)=0$ by the preceding two-factor argument, with the
other factors contributing no invariant vectors.  Composing with
an inner automorphism gives the corresponding statement for
automorphisms over an inner automorphism of $Q$.

\paragraph{The exact symmetry maps and the kernel family.}
The program uses the actual $S_5$ conjugations in each component
map.  For such a conjugation $a$, it finds $L_a$ and $b_a$ with
\[
 c(a g)=L_a c(g)+(\rho(a g)-1)b_a                       \tag{Aux.3}
\]
for all sixty $g$.  Here $L_a$ is a nonzero scalar multiple of
the adjoint matrix of the specified $\operatorname{PGL}_2(\F_5)$
representative.  The scalar and vector are unique and retained
in the output.  On a target summand $(i,j)$ with source summand
$(s_i,s_j)$, put
\[
 A=L_i c(q_{s_i}),\quad D=L_j c(q_{s_j}),\quad
 r=\rho(a_jq_{s_j})b_j,
 \qquad
 J_a(q)_{i,j}=-b_i\otimes D+A\otimes r-b_i\otimes(r-b_j).
\]
With the corresponding block permutation map $L_a$ on $M$,
the formula
\[
 \Phi_a(q,m)=(a(q),L_a m+J_a(q))                         \tag{Aux.4}
\]
is an automorphism of (Aux.2).  Substituting (Aux.3) proves
\[
 L_a B(q,r)+J_a(qr)=J_a(q)+\rho(aq)J_a(r)+B(aq,ar),
\]
which is exactly its multiplication condition.  The program
also checks this equality on all forty-nine pairs of the seven
named elements $1,X,Y,Z,X^{-1},Y^{-1},Z^{-1}$, for each of
$a=\theta,\omega$.  Every other lift of the same component map
is obtained by adding $(1-\rho(aq))b$ to (Aux.4).

Here $X,Y$ denote the specified base images of $x,y$ and
$Z=(XY)^{-1}$.  An epimorphism $F\to L$ over the fixed map to
$Q$ is specified by arbitrary module coordinates $X_v,Y_v$.
It is surjective by $M\subset\Phi(L)$.
Two such maps have the same kernel exactly when they differ
by an automorphism over $Q$, hence by simultaneous conjugation
by $M$.  The space of kernels is therefore the affine space
\[
 \F_5^{108}/
   \{((1-\rho(X))m,(1-\rho(Y))m):m\in M\},
                                                        \tag{Aux.5}
\]
of dimension $54$.  The denominator has dimension $54$
because $M^Q=0$.

The exact first linear system imposes, on the lifts,
\[
 \Theta(X)=Y,\quad\Theta(Y)=X,\quad
 \Omega(X)=Y,\quad\Omega(Y)=(XY)^{-1}.                  \tag{Aux.6}
\]
Its unknowns are $X_v,Y_v,b_\theta,b_\omega\in M$;
$b_\theta,b_\omega$ are the inner corrections in (Aux.4).
Writing $R_W=\rho(W)$ and $L_\theta,L_\omega$ for the
two block maps, its matrix before choosing the conjugation
gauge is
\[
 \begin{pmatrix}
 L_\theta&-1&1-R_Y&0\\
 -1&L_\theta&1-R_X&0\\
 L_\omega&-1&0&1-R_Y\\
 R_Z&L_\omega+R_ZR_X&0&1-R_Z
 \end{pmatrix}.                                        \tag{Aux.7}
\]
The right side is the column with entries
$-J_\theta(X)$, $-J_\theta(Y)$, $-J_\omega(X)$, and
$-J_\omega(Y)-R_Z(B(X,Y)+B(XY,Z))$.
The program chooses $54$ independent coordinate rows of
the denominator in (Aux.5) and sets those coordinates to zero.
Row reduction gives an affine solution space of dimension
$9$, parametrizing exactly $5^9$ kernels satisfying (Aux.6).

This family supplies a necessary finite test for an actual
$GT_0$ involution.  Precomposition on epimorphisms, followed
by a specified lift on the target, acts affinely on (Aux.5).
Changing that target lift by inner conjugation does not change
the kernel.  Inner automorphisms of $F$ also act trivially on
kernels.  For an actual involution, the outer subgroup generated
by it and the geometric $S_3$ therefore acts through
$C_2\times S_3$ on this affine space.  Averaging the orbit of
a point, with division by $12$ in $\F_5$, gives a common fixed
kernel.  The induced involution of $L$ is inner by the
automorphism calculation above.  This explains the scope of
the simultaneous system without assuming that one arbitrarily
chosen epimorphism kernel is invariant.

\paragraph{The involution and peripheral equations.}
For every base tuple $t$ surviving the D.38 triangle conditions
on $I$, introduce $T=(t,m)$ and impose
\[
 T^2=1,\qquad TXT^{-1}=X^{-1},\qquad
                  \Omega(T)=Y^{-1}\Theta(T).            \tag{Aux.8}
\]
These are affine linear equations over $\F_5$, jointly in
the nine parameters of (Aux.7) and the $54$ coordinates of
$m$.  In particular the third equation is the full finite
three-cycle compatibility, rather than an independent local
inversion test.  To specify the equations completely, their
left sides before inserting the nine-parameter solution are
\[
 \begin{split}
 &(1+R_t)m,\\
 &(1-R_{X^{-1}})m+(R_t+R_{X^{-1}})X_v,\\
 &(L_\omega-R_{Y^{-1}}L_\theta)m
 +(1-R_{\omega t})b_\omega
 -R_{Y^{-1}}(1-R_{\theta t})b_\theta+R_{Y^{-1}}Y_v.
 \end{split}                                           \tag{Aux.9}
\]
Their right sides, in the same order, are
\[
 \begin{split}
 &-B(t,t),\\
 &-B(t,X)-R_{X^{-1}}B(X,X^{-1})+B(X^{-1},t),\\
 &-J_\omega(t)+R_{Y^{-1}}J_\theta(t)
       -R_{Y^{-1}}B(Y,Y^{-1})+B(Y^{-1},\theta t).
 \end{split}
\]
All matrices and constants are specified by (Aux.1)--(Aux.4)
and the retained D.38 factor maps.

The five independent component sets are
$\{3,5\},\{6,7\},\{8\},\{11,14\},\{13,15\}$,
with respectively $3,2,2,5,5$ choices.  Thus exactly
$300$ base tuples occur.  The executed outcome is
\[
 \begin{array}{l|r}
 \text{quantity}&\text{value}\\ \hline
 \text{base tuples in the restricted D.38 conjugacy model}&150\\
 \text{base tuples outside that model}&150\\
 \text{base tuples admitting (Aux.6) and (Aux.8)}&300\\
 \text{solution dimension over }\F_5\text{ for every tuple}&6\\
 \text{dimension of its projection to the kernel family}&3\\
 \text{compatible kernels for each base tuple}&125\\
 \text{involutions over each compatible kernel}&125.
 \end{array}                                           \tag{Aux.10}
\]
The last count uses center-freeness: different elements $T$
give different inner automorphisms.  The total is
$300\cdot5^6=4\,687\,500$ kernel--involution pairs,
split equally between the base tuples inside and outside
the restricted D.38 conjugacy model.

For a particular solution of each of the $300$ systems the
program also evaluates the original group multiplications,
independently of the row-reduction equations.  With
$f=\Theta(T)T^{-1}$ it verifies
\[
 \Theta(f)f=1,\qquad f\operatorname{Inn}(T)(f)=1,
 \qquad
 \Omega^2(f)Z^{-1}\Omega(f)Y^{-1}fX^{-1}=1.              \tag{Aux.11}
\]
These identities hold throughout the affine solution sets:
(Aux.6) gives the actual word maps on the generating pair,
and (Aux.8) implies (Aux.11) by the cancellations used in
D.38.  Thus this refinement tests the simultaneous finite
peripheral and power equations and still eliminates no base
tuple.  It does not establish a solution of D.35's conjugator
equation, or a compatible lift through all finite quotients.

The executed programs are
\path{code/part_a/a5_cross_module.py} and
\path{code/part_a/a5_cross_peripheral_lifts.py}.
Their adjacent JSON files retain the cocycle, the component
orbit, the cohomology corrections, every one of the $300$
base tuples, and all dimensions in (Aux.10); adjacent text
files retain their standard output.
